\begingroup
\sloppy
\setlength{\emergencystretch}{3em}
\raggedright
\section{Unified Mobile Pilot Guardian}
\label{sec:pilot-unified-mobile-guardian}

Pilot Guardian is now available inside the unified Personal phone application as a priority-help
surface.  It reuses the existing Guardian authority and remains deliberately separate from ordinary
messages, calls and contact permissions.  Its present production claim is narrow: it can prepare
and attempt an urgent alert to a contact explicitly authorized by the protected person or their
guardian.  It does not diagnose a fall, contact an emergency service or claim ambulance dispatch.

\subsection{Separate Authority and Protected Person}

The mother issues three distinct capabilities: \texttt{guardian.read},
\texttt{guardian.configure} and \texttt{guardian.alert}.  Requests are signed by the paired phone
and bound to its active certificate, lease and private identity.  A phone may protect its own person
or a child profile explicitly linked to that adult as guardian; it cannot open or act upon an
unrelated adult's Guardian state.

An emergency contact permission names one existing private contact, one delivery channel, whether
emergency interruption is allowed and whether location may be shared.  The permission explicitly
states that it grants no ordinary-contact authority.  Location sharing additionally requires a
separate permission receipt and is omitted from alerts when that proof is absent.

\subsection{Confirm-before-send Incident Flow}

Pressing the large help control creates an incident in
\texttt{awaiting\_large\_confirmation}.  No alert is sent at this point.  The surface presents the
number of authorized routes and separate controls to confirm or cancel a false alarm.  Confirmation
locks the exact permission identifiers and then attempts their registered providers.  A result is
\texttt{alerts\_verified} only when a provider returns both a positive verification and provider
alert identifier.  Missing, failed or unverified adapters produce
\texttt{connectivity\_fallback\_required} and a prominent instruction to use local emergency
services directly.

Create, configure, confirm, deliver and cancel paths are idempotent.  Network retransmission of an
identical signed operation returns its recorded result; reuse of the key for changed details is
rejected.  In particular, a repeated phone response cannot silently produce a second provider
alert.

\subsection{Safety State and Verification}

The interface always displays that emergency-service dispatch is not connected.  Each returned
state includes \texttt{ambulance\_dispatched=false}, an explicit delivery-verification flag and a
fallback flag.  Sensor observations remain prompts rather than diagnoses or automatic dispatch
authority.

Focused Guardian, identity-boundary and mobile-surface tests verify separate permissions, the
large-confirmation boundary, fail-closed delivery, honest dispatch claims and replay protection.
The optimized Next.js production build compiles successfully.  Automated verification registered
no real contact, sent no alert and contacted no emergency service.

\endgroup
